\section{Addressing}

An infinite space is useless to a computer if you cannot say where in it you are. The base settles this completely. Every dock carries a short unique name, and stepping to any neighbor is fixed arithmetic on that name, with no lookup table and no search.

\begin{result}[A solved address for every dock]
Every dock has a unique finite address whose length grows like the logarithm of its distance, and each of its twenty-four neighbors is reached by a closed-form operation on that address. Greedy routing by address alone delivers every message along a path a fraction longer than the true shortest, verified out to two million docks. \cite{pollard2026vibetest}
\end{result}

The idea is most transparent in Margenstern's classic two-dimensional case, the heptagrid $\{7,3\}$, where the growth is the Fibonacci recurrence and the addresses are \emph{Zeckendorf} strings \cite{margenstern2008, margenstern2013}. Number the cells ring by ring, then rewrite each number in the Fibonacci base, whose place values are $1, 2, 3, 5, 8, 13, \ldots$. By Zeckendorf's theorem every integer is a unique sum of non-adjacent Fibonacci numbers, so no address ever has two ones in a row, which makes the valid addresses a regular language that a tiny automaton recognizes.

\begin{table}[ht]
\centering
\begin{tabular}{@{}rll@{}}
\toprule
cell & Fibonacci sum & address \\
\midrule
1  & 1         & 1 \\
2  & 2         & 10 \\
3  & 3         & 100 \\
4  & $3 + 1$     & 101 \\
5  & 5         & 1000 \\
6  & $5 + 1$     & 1001 \\
7  & $5 + 2$     & 1010 \\
8  & 8         & 10000 \\
9  & $8 + 1$     & 10001 \\
10 & $8 + 2$     & 10010 \\
11 & $8 + 3$     & 10100 \\
12 & $8 + 3 + 1$ & 10101 \\
\bottomrule
\end{tabular}
\caption{Zeckendorf addressing on the heptagrid $\{7,3\}$.}
\end{table}

Routing is pure string surgery. The parent is found by stripping a digit, a child by extending by a legal digit. To route from $A$ to $B$ you compare addresses from the left, walk up to the common ancestor, then walk down. Sending a message from cell 12, with address 10101, to cell 9, with address 10001, the shared prefix is 10, which names cell 2 as the common ancestor. The path goes up from 12 to 7 to 4 to 2, then down from 2 to 3 to 5 to 9. The full route is $12, 7, 4, 2, 3, 5, 9$, six hops, with every cell comparing only two short strings and no map consulted. Heptagrid routing delivered 100 percent at low stretch in the runs. \cite{pollard2026vibetest} This is the clean two-dimensional template. The $\{3,4,3,4\}$ case carries the same idea, with route-from-the-root, parent-by-strip, and common-prefix routing, but with a richer numeration whose growth root is near $18.3$ rather than the golden ratio, so the base is the twenty-four-way child index rather than Zeckendorf binary, and with no cousins.

A worked example makes it concrete. Addresses on $\{3,4,3,4\}$ are dot-separated digits. Each digit is the index of a cell among its parent's children, ordered deterministically by embedding, and the whole string is the route from the root. The root is the empty address written $e$. Shell-1 cells are single digits, shell-2 cells are two digits, and child counts vary by son type. The following cells are real output of the addressing code. \cite{pollard2026vibetest}

\begin{table}[ht]
\centering
\begin{tabular}{@{}rrlrr@{}}
\toprule
cell & shell & address & parent & children (son type) \\
\midrule
0  & 0 & $e$   & --- & 24 \\
1  & 1 & 0   & 0 & 23 \\
2  & 1 & 1   & 0 & 22 \\
3  & 1 & 2   & 0 & 21 \\
\dots & 1 & \dots & 0 & \dots \\
24 & 1 & 23  & 0 & 15 \\
25 & 2 & 0.0 & 1 & 22 \\
26 & 2 & 0.1 & 1 & 21 \\
\bottomrule
\end{tabular}
\caption{A few $\{3,4,3,4\}$ cells with shell, address, parent, and child count.}
\end{table}

A routing walk from cell $A$ to cell $B$ is pure address arithmetic. Compare the two addresses from the left, walk up from $A$ to the common ancestor by dropping the last digit to reach the parent, then walk down to $B$ by appending $B$'s remaining digits to reach the named child. Take $A$ to be cell 481, address 0.0.0, and $B$ to be cell 503, address 0.1.0. Their longest shared prefix is 0, so the common ancestor is cell 1, address 0.

\begin{table}[ht]
\centering
\begin{tabular}{@{}rrll@{}}
\toprule
step & cell & address & operation \\
\midrule
0 & 481 & 0.0.0 & start ($A$) \\
1 & 25  & 0.0   & drop last digit, go to parent \\
2 & 1   & 0     & drop last digit, reach the common ancestor \\
3 & 26  & 0.1   & append digit 1, go to child \\
4 & 503 & 0.1.0 & append digit 0, reach $B$ \\
\bottomrule
\end{tabular}
\caption{A routing walk from cell 481 to cell 503, four hops.}
\end{table}

Four hops, which equals the shortest-path distance, is exactly the address arithmetic the routing uses, with no map and no coordinates. The neighbour set of $A$ is likewise reconstructed from its name alone. It is parent 25, plus 21 children, plus 2 confluence alternate-parents, plus 0 alternate-children, totalling $A$'s full 24 facet-neighbours.

What makes this exact at any depth is that the address is really a word in the mirror reflections, and reflection groups have a solved word problem. Naming a dock by its rounded position works to moderate depth, then fails, because in hyperbolic space the cells crowd exponentially toward the boundary and floating-point error builds along a long reflection path until distinct docks collide onto one key. The cure is to name each dock by the shortest reflection word that reaches it, ties broken alphabetically, a normal form with exactly one word per dock. Two facts make it work. Two words name the same dock precisely when one turns into the other by braid moves alone, and every word shortens to a reduced one by cancellation, so equality is decidable in finite steps. And every reflection group is automatic, its normal-form words a regular language read by a finite machine, so one small automaton enumerates the entire infinite crystal, one word per dock, with no map, no crowding, and no drift. This is the same Fibonacci idea generalized from the flat two-dimensional case to every honeycomb, including the three- and four-dimensional ones where a hand-built scheme does not reach.

This matters twice over. For anyone simulating the substrate, any region can be addressed and stepped exactly, with no approximation and no stored map. For the theory itself, a pattern living in the mesh, a self, has a place it can name and return to, a coordinate on every dock within reach. The tree of the mesh is not only a shape. It is a memory with an address on every node.
